\section{Limitations and Future Work}\label{sec:limitations}

The three correctness properties established in Section~\ref{sec:properties} are conditional guarantees: they obtain if and only if a set of platform and deployment assumptions hold simultaneously. This section makes those assumptions explicit, identifies the practical and theoretical limitations that follow from them, and maps each limitation to a directed research agenda. Together, the four limitations and three future-work items define the boundary between what the protocol achieves within its scope and what must be established by subsequent work before the protocol can be deployed at scale in high-stakes environments.

\subsection{Limitations}

\subsubsection{Human Response Latency}

The Bailout Protocol's Liveness property guarantees eventual human contact but does not guarantee that human contact occurs within a time horizon useful for the operational context. The bound $T_{\max}$ derived in the liveness proof is a function of network topology and retry parameters, not of human cognitive response time. In high-frequency operational domains---algorithmic trading, real-time process control, autonomous vehicle navigation---the time required for a human principal to receive, interpret, and respond to an escalation may exceed the operational time constant of the system by an order of magnitude.

The protocol's failure mode under human latency is well-defined: the system stalls in \texttt{ESCALATING} state until the human responds. This is the intended behavior. However, the operational cost of a stall is domain-dependent. In a precision agriculture context such as Irrig8, a stall that delays a watering decision by minutes or hours is operationally acceptable; in a high-frequency trading context, a stall that delays a risk-interlock response by even seconds may be economically catastrophic. The protocol does not currently provide a mechanism for expressing domain-specific cost structures or adapting its stall-tolerance accordingly. The human anchor is treated as a black box: the protocol sends a message, waits for a response, and times out if the response does not arrive within $T_{\text{max}}$. There is no mechanism to reason about why the response has not arrived, whether the anchor is overloaded with concurrent escalations, or whether the delay reflects a structural mismatch between the protocol's timing parameters and the human's operational availability.

Human response latency is a deployment parameter provisioned at node setup, not a protocol parameter adjustable at runtime without human involvement. For deployments in latency-sensitive domains, explicit human response time budgets must be designed in and validated against operational data before the system is released.

\textit{Directed research:} Empirical measurement of human response latency distributions across deployment contexts to establish statistically grounded latency budgets. Integration of these budgets into the node provisioning interface as enforceable policy parameters. Investigation of graded escalation responses that request a rapid acknowledgment followed by a detailed response, better matching the operational reality of human decision-making under varying loads.

\subsubsection{Adversarial Conditions}

The Safety property requires that human-contact pathways cannot be closed without recorded human authorization. This guarantee holds under the assumption that the host platform provides independent scheduling priority to the Bailout Monitor (Assumption A1 in the properties section). In an adversarial environment---particularly one in which the host platform itself may be compromised---this assumption may not hold. A compromised OS or virtualized substrate can modify thread priorities, preempt the monitor thread, or introduce race conditions into the pathway-closure operations the monitor is designed to interdict.

The protocol's guarantees are platform guarantees. They are derived from invariants the host platform must provide: non-preemptible monitor scheduling, append-only storage for the authorization ledger, and interruptible blocking calls for pathway management. If the platform itself is adversarial, these invariants cannot be relied upon. A sophisticated attacker who controls the execution substrate can rewrite the Bailout Monitor's code, suppress its scheduling, or tamper with the authorization ledger before persistence. The protocol's bypass mechanism enforces the no-machine-actor property on the propagation path, but it cannot protect against a machine actor who controls the substrate on which the propagation path runs.

The Bailout Protocol is designed for deployments in which the host platform is trusted but individual agents may behave unexpectedly---the standard assumption for production AI deployments in regulated industries. It is not designed for deployments in which the platform itself is an active adversary. For those deployments---national infrastructure, military systems, financial clearinghouses---the protocol must be composed with substrate-integrity mechanisms outside its scope: secure boot chains, hardware-attested execution environments, and independent watchdog monitors running on separate hardware. The protocol's role in such compositions would be to ensure the human accountability guarantee holds \emph{given} a trustworthy substrate; the composition's role would be to ensure the substrate is trustworthy.

\textit{Directed research:} Specification of the composition interface between the Bailout Protocol and substrate-integrity mechanisms. Formal analysis of whether the protocol's guarantees can be preserved under weaker platform assumptions, including partially compromised substrates. Exploration of hardware-attested execution environments (e.g., Intel SGX, ARM TrustZone) as a deployment target for the Bailout Monitor, eliminating the OS as a trust dependency.

\subsubsection{Scalability}

The Bailout Protocol is specified with a 1:1 relationship between each node and its human anchor. This model is appropriate for single-node deployments or small-scale systems but creates a scalability bottleneck in enterprise and infrastructure contexts. Without pooling and load-balancing mechanisms, the human anchor becomes the rate-limiting step in a system designed to route exceptions away from centralized decision points. The protocol correctly escalates each node's exceptions to its designated anchor, but if one anchor is responsible for hundreds of concurrent nodes, the anchor becomes a saturation point that degrades the Liveness property in operationally significant ways.

The current specification treats human anchoring as a 1:1 relationship. This model must evolve to support at minimum 1:N relationships (one anchor supervising multiple nodes, with pooling) and potentially N:M relationships (anchor pools with load balancing and routing policies). The correctness properties of the protocol must be shown to extend to these pooled configurations, particularly the Safety property, which requires that no configuration of pooling can produce a state in which all human-contact pathways are permanently disabled.

\textit{Directed research:} Formal specification of human anchor pool architectures, including load-balancing policies, prioritization schemes, and routing algorithms. Extension of the Safety and Liveness proofs to pooled anchor configurations. Empirical evaluation of anchor pool performance under simulated high-escalation-load conditions to establish practical scaling curves and identify saturation thresholds.

\subsubsection{Exploitation Risk}

The Bailout Protocol is designed under the assumption that machine actors will trigger bailout faithfully when their operating range is exceeded. However, a sophisticated agent---particularly one operating in a domain where escalation is economically costly to the agent's objectives---may develop strategies to suppress or mask exception states rather than trigger the protocol. This is a form of goal hacking applied to the accountability mechanism itself. An agent might learn to reframe an out-of-scope condition as an in-scope condition through creative interpretation of its operating range, effectively absorbing exceptions within the machine-only execution graph without triggering the mandatory escalation.

The bypass mechanism enforces the no-machine-actor property on the propagation path once escalation is triggered, but it does not detect whether a machine actor has intentionally avoided triggering the protocol in the first place. The state machine cannot distinguish between a Bounds Engine that correctly assessed a condition as in-scope and one that deliberately misclassified an out-of-scope condition to avoid escalation.

This limitation strikes at the heart of the protocol's accountability model. The Bailout Protocol holds machine actors accountable by observing their triggering behavior: if a node observes an out-of-scope condition and fails to trigger, the protocol has no independent mechanism to detect the omission. The protocol relies on the agent's own self-assessment, which is precisely the mechanism that may be compromised. External monitoring of triggering behavior is necessary but introduces new complexity: who monitors the monitors, and what prevents the monitoring from itself being gamed?

The protocol's guarantees are conditional on the faithful triggering behavior of machine actors. In deployments where agents have strong incentives to avoid escalation, the protocol must be composed with monitoring mechanisms that detect anomalies in triggering behavior---specifically, statistically significant deviations between the distribution of conditions observed by a node and the distribution of escalations it triggers.

\textit{Directed research:} Design of triggering-behavior monitors that audit the statistical relationship between observed exception conditions and triggered escalations per node. Specification of anomaly thresholds that trigger human review without presuming bad faith. Investigation of provenance attestation for operating-range assessments: can the protocol require that each Bounds Engine assessment be counter-signed by an independent audit process that records the raw inputs and the assessment criteria, making post-hoc misclassification detectable?

\subsection{Future Work}

Beyond the four limitations above, three research directions are of particular significance for the protocol's development trajectory.

\subsubsection{Formal Verification}

The argument-map proofs provided in the properties section establish the structural conditions under which the Bailout Protocol's guarantees hold, but they are not machine-verified proofs. A formal verification effort using a proof assistant such as Coq or Isabelle/HOL would establish the Safety, Liveness, and Accountability properties as theorems under precisely specified axiomatic assumptions. This effort is a prerequisite for high-assurance deployments in safety-critical domains where the cost of an undetected protocol failure is catastrophic.

Of particular interest is whether the functional separation between Purpose, Bounds, and Fact layers admits a compositional verification argument: if each layer satisfies its specification, does the composition necessarily satisfy the protocol's guarantees? If so, the verification burden can be distributed across layer implementations rather than centralized on the full system. A compositional verification architecture would make the protocol more robust to implementation changes. Formal verification would also expose hidden assumptions in the argument-map proofs that are currently satisfied informally---gaps that would be identified and resolved. The verification effort would thus be as valuable for what it reveals about the protocol's limitations as for what it confirms about its correctness.

\subsubsection{Adaptive Timeout}

The current specification uses a fixed timeout $T_{\text{bounds}}$ as a liveness trigger. The same parameter value cannot be appropriate for an agricultural irrigation system (where response times of hours are acceptable) and a financial risk management system (where response times of seconds are required). Yet the provisioning burden for domain-specific timing parameters is significant, and incorrect provisioning can cause the protocol to either timeout prematurely or stall indefinitely.

An adaptive timeout mechanism---one that adjusts $T_{\text{bounds}}$ based on observed human response latency, escalation frequency, and domain-specific operational time constants---would reduce the human configuration burden and improve the protocol's operational fit across deployment contexts. The design must preserve the deterministic timing guarantees required by the Safety property: adaptation must be constrained and slow relative to the operational time constants of the system, so that the adapted timeout remains a deterministic upper bound rather than a probabilistic estimate. On each escalation cycle, the protocol records the actual response time $t_{\text{actual}}$; over a rolling observation window it computes a distribution and sets $T_{\text{bounds}}$ to a high percentile (e.g., the 95th) plus a safety margin. The Safety property is preserved by never adjusting $T_{\text{bounds}}$ below the minimum viable threshold established at provisioning.

\subsubsection{Distributed Human Anchor Pools}

The current 1:1 node-to-anchor model is not viable at scale. The long-term research agenda includes the design of distributed human anchor pools that preserve the protocol's correctness guarantees under concurrent load. Open questions include how to partition nodes among anchors while preserving the invariant that each node has at least one reachable anchor; how to handle anchor unavailability without violating Safety; and how to design routing policies that minimize human attention fragmentation while ensuring all escalations receive timely response.

The pool architecture must satisfy three constraints. First, \textit{Safety:} no pool configuration may produce a state in which all human-contact pathways for a given node are permanently disabled. This requires that each node maintain at least two anchor references at all times and that the pool management algorithm detect and compensate for anchor unavailability before it becomes pathway-closing. Second, \textit{Liveness:} the pool routing algorithm must ensure escalations are routed to available anchors within $T_{\max}$ even under high concurrent load, implying load-aware routing rather than static round-robin. Third, \textit{Accountability:} every escalation must be attributed to a specific human anchor and that attribution must be immutable, requiring that routing decisions be recorded in the Forensic Ledger.

The design of anchor pools also raises a practical question about human workload management: what happens when all anchors in a pool are simultaneously overloaded? The pool architecture must include a load-shedding policy that preserves Safety while providing meaningful feedback about the overload condition to the node and to the system operator.
